% What we need from Chapters 3 and 4a --- spec \S5.1.
% Boxed formulas of Project_specification.md \S6 with local labels; all
% later sections reference these. Full derivations: Chapters 3 and 4a.

\section{What we need from Chapters 3 and 4a}
\label{sec:needed}

This chapter embeds the single-cell theory of Chapters~3 and~4a in a
population. Nothing new from the single cell is required --- but several of
its closed forms will be quoted constantly, and they are collected here once,
boxed, with the labels the rest of the chapter uses. Chapter~3 supplies the
derivations; Chapter~4a the distribution theory; this section is the
reference card.

Throughout: $\lambda$, $\mu$, $\delta$ are the intracellular birth, death and
catastrophe rates per replicative unit; $H$ is the absorbing rupture state;
$\xt$ the intracellular load, $\wt$ the released count, $\tau$ the burst
time, and $\Kb$ the burst size. We write $w=\ee^{\theta t}$, and one
naming decision is made now and kept: the burst-time density, which is
$\delta J(t)$ in both earlier chapters, is written as $\delta J$ in this
chapter, because $\varphi$ is reserved for the partial-release fraction of
\S\ref{sec:boolean}.

\subsection{Roots and identities}
\label{sec:needed-roots}

The roots of the quadratic $\lambda f^2-(\lambda+\mu+\delta)f+\mu$ are
\begin{equation}
a=\eta+\sqrt{\eta^2-\frac{\mu}{\lambda}},\qquad
b=\eta-\sqrt{\eta^2-\frac{\mu}{\lambda}},\qquad
\eta=\frac{\lambda+\mu+\delta}{2\lambda},
\end{equation}
with $b<1<a$, and the derived constants are
\begin{equation}
A=a-1,\qquad B=1-b,\qquad \theta=\lambda(a-b),\qquad \kappa=1+\frac{\delta}{2\lambda}.
\end{equation}
The identities that simplify everything downstream are
\begin{equation}
\boxed{ab=\frac{\mu}{\lambda},\quad a+b=\frac{\lambda+\mu+\delta}{\lambda},\quad
AB=\frac{\delta}{\lambda},\quad A+B=a-b,\quad \theta=\lambda(a-b).}
\label{eq:AB}
\end{equation}

\subsection{The three fixation functions}

The no-burst probability, the internal-extinction probability, and the
non-fixation probability are
\begin{equation}
\boxed{I(t)=\frac{aB+bAw}{B+Aw},\qquad
D(t)=\frac{ab\,(w-1)}{aw-b},\qquad
\Ifix(t)=I(t)-D(t)=\frac{(a-b)^2\,w}{(B+Aw)(aw-b)}.}
\label{eq:IDIhat}
\end{equation}
$I(\infty)=D(\infty)=b$ --- $b$ is the probability that the cell clears
internally without ever bursting --- while $\Ifix(0)=1$,
$\Ifix'(0)=-(\mu+\delta)$ and $\Ifix(\infty)=0$. The population survival
kernel of \S\ref{sec:renewal} is $\Ifix$ itself.

\subsection{Moments and mean release}

With $I'=-\delta J$ and $V'=\delta K$,
\begin{equation}
\boxed{J(t)=\frac{(a-b)^2\,w}{(B+Aw)^2}=-\frac{\lambda}{\delta}(I-a)(I-b),}
\label{eq:J}
\end{equation}
\begin{equation}
\boxed{K(t)=\Bigl[1+\frac{2\lambda}{\delta}(1-I)\Bigr]J
=\frac{2\lambda}{\delta}(\kappa-I)J,}
\label{eq:K}
\end{equation}
and the cumulative mean release, a quadratic in $1-I$ alone,
\begin{equation}
\boxed{V(t)=(1-I)\Bigl[1+\frac{\lambda}{\delta}(1-I)\Bigr].}
\label{eq:V}
\end{equation}
Its late-time value, the total expected output of one cell, is
\begin{equation}
\boxed{V_\infty=\frac{a(1-b)}{a-1}=\frac{aB}{A}
=\frac{\lambda-\mu}{\delta}(1-b)+1,
\qquad
\mean{\Kb\mid\text{burst}}=\frac{V_\infty}{1-b}=\frac{a}{a-1}.}
\label{eq:Vinf}
\end{equation}
Note again that $V_\infty$ counts the internally cleared cells with their
release of zero; the conditional mean removes them.

\subsection{Distribution-level facts}

Three results of Chapter~4a are used as black boxes here. First, the
burst-time density is
\begin{equation}
-I'(t)=\delta J(t),
\label{eq:phi}
\end{equation}
defective with total mass $1-b$. Second, the burst size, unconditional,
obeys
\begin{equation}
\pr{\Kb=k}=\frac{\delta}{\lambda}\,a^{-k},\qquad k\ge1,
\label{eq:burstlaw}
\end{equation}
with the mass $b$ at ``no burst''; conditioned on bursting this becomes
$(a-1)a^{-k}$, which is exactly the quasi-stationary distribution of the
load (both are geometric on $\{1,2,\ldots\}$ with ratio $1/a$). Third,
given that the burst occurs at time $t$, its size is a size-biased draw
from the state probabilities:
\begin{equation}
\boxed{\mean{\Kb\mid\tau=t}=\frac{K(t)}{J(t)}
=1+\frac{2\lambda}{\delta}\bigl(1-I(t)\bigr),}
\label{eq:sizebias}
\end{equation}
tending at late time to $(a+1)/(a-1)$, the size-biased quasi-stationary
mean.

\subsection{The second moment of the release}

The whole-host variance prediction of \S\ref{sec:open-problems} starts from
\begin{equation}
\boxed{\mean{\wt^2}
=\frac{2(\lambda-\mu)}{\delta}V
-\frac{\lambda+\mu}{\delta}I
-K
+\frac{\lambda+\mu}{\delta}
+1,}
\label{eq:EW2}
\end{equation}
and $\mathrm{Var}(\wt)=\mean{\wt^2}-V^2$.

\subsection{The mean productive lifetime}

The integral of the survival kernel, used in \S\ref{sec:peff-section} to
anchor the effective parameters, is
\begin{equation}
\boxed{\mean{T_{\rm prod}}=\int_0^\infty\Ifix(t)\,\dd t
=\frac{1}{\lambda}\log\frac{a}{a-1}
=\frac{1}{\lambda}\log\mean{\Kb\mid\text{burst}}.}
\label{prop:lifetime}
\end{equation}

\subsection{Multiplicity of infection}

A cell founded by $k$ particles releases superlinearly more. Its moments
and release flux are
\begin{equation}
J_k(t)=kJ(t)I(t)^{k-1},\qquad
K_k(t)=kK(t)I(t)^{k-1}+k(k-1)J(t)^2I(t)^{k-2},
\end{equation}
\begin{equation}
\boxed{g_k(t)=\delta K_k(t)
=\delta\lrb{kK I^{k-1}+k(k-1)J^2I^{k-2}},}
\label{eq:gk}
\end{equation}
and its lifetime yield, for $k\ge2$,
\begin{equation}
\boxed{V_\infty^{(k)}
=(1-b^k)
+\frac{2\lambda}{\delta}\Biggl[\frac{1}{k+1}-b^k+\frac{k\,b^{k+1}}{k+1}\Biggr]
+k(k-1)\frac{\lambda}{\delta}\int_1^b (x-a)(x-b)\,x^{k-2}\,\dd x,}
\label{eq:Vk}
\end{equation}
simplifying to $V_\infty^{(k)}=k+\lambda/\delta$ when $\mu=0$. The
conditional mean burst $V_\infty^{(k)}/(1-b^k)$ increases with $k$:
multiply infected cells burst more often, and bigger.

\subsection*{Scope}

Everything in this section is quoted; the derivations live in Chapters~3
and~4a, and the verification of the population-level consequences of what
follows is recorded in Appendix~\ref{app:verification}. The population
notation of this chapter --- $\Icell$, $\Vfree$, $\gamma$, $c$, $T$,
$p$, $d_{\Icell}$ --- is chosen to avoid collision with all of the above,
and is collected in Table~\ref{tab:notation-pop}.

\begin{table}[H]
\centering
\renewcommand{\arraystretch}{1.25}
\small
\begin{tabular}{@{}llp{0.44\textwidth}@{}}
\toprule
Symbol & Meaning & Notes \\
\midrule
$T$ & target-cell count & held fixed \\
$\gamma$ & infection rate & \\
$\Icell,\Vfree$ & infected cells, free particles & population counts \\
$c$ & virion clearance & \\
$d_{\Icell}$ & infected-cell removal (phenomenological) & not $\delta$ \\
$p$, $p_{\rm eff}(r)$ & constant / effective release rate & \\
$g(a)$ & release kernel, $\delta K(a)$ & expected flux at age $a$ \\
$S(a)$ & survival kernel, $\Ifix(a)$ & \\
$q$ & $\gamma T/(\gamma T+c)$, infection success & branching \\
$m$ & $qV_\infty$, branching mean offspring & \\
$L$ & $a(1-b)=a-\mu/\lambda$, flooding parameter & \\
$\varphi$ & boolean partial-release fraction (\S\ref{sec:boolean}) & \\
$\zeta$ & intensity decay rate of Model~10 (\S\ref{sec:self-excite}) & not $\kappa$ \\
\bottomrule
\end{tabular}
\caption{Population-level notation of this chapter.}
\label{tab:notation-pop}
\end{table}
